% SOURCE_AUTHORITY: sga5_fr_workpass.tex, lines 11209--11229 (LF-normalized unit).
% SOURCE_SHA256: 791F4EFFC5E02832D5D77ED03518C8156D6F07E4C8238B03545DB93D883FBB28
Como $X\times0$ é linearmente equivalente a $X\times1$ em $X\times\mathbb P^1$, temos a igualdade $f_1^*u_*(1)=f_1^*(X\times1)$, de onde, como $\widehat N\cap f_1^{-1}(X\times1)=\emptyset$,
\[
(j_1)_*g_1^*(a)\cdot f_1^*u_*(1)=0.
\tag{9.4.6}
\]
Comparando com (9.4.5), deduz-se que
\[
(j_1)_*g_1^*(b_0)=-v_*(1)\bigl[(j_1)_*g_1^*(a)+(j_1)_*(z)\bigr].
\tag{9.4.7}
\]
Por outro lado, é claro que $W'\cap X'=\emptyset$; de fato, quando muito poderiam cortar-se sobre $Y$, onde um dá a seção zero e o outro o hiperplano no infinito de $\widehat N$. Por conseguinte, $\beta^*v_*(1)=0$, de onde, por (9.4.7),
\[
\beta^*(j_1)_*g_1^*(b_0)=0.
\tag{9.4.8}
\]
Mas, como já foi observado (cf. demonstração de 9.3.2), $\beta$ e $j_1$ são transversais, de modo que (9.4.8) proporciona
\[
0=\delta_*s^*g_1^*(b_0)=\delta_*(b_0).
\tag{9.4.9}
\]
Finalmente, aplicando $(f_2)_*$ a (9.4.9), obtém-se $\gamma_*(b_0)=0$, de onde $b_0=0$, já que $\gamma_*$ é injetivo ($\gamma$ admite uma retração evidente).
